\documentclass[11pt]{amsart}

\usepackage[T1]{fontenc}
\usepackage{lmodern}
\usepackage{microtype}
\usepackage{mathtools,amssymb,amsthm}
\usepackage[hidelinks]{hyperref}

\hypersetup{
  pdftitle={A union of two non-DDMO copies of K(3,3) is difference distance magic orientable}
}

\newtheorem{theorem}{Theorem}
\newtheorem{lemma}[theorem]{Lemma}
\newtheorem{proposition}[theorem]{Proposition}
\newtheorem{corollary}[theorem]{Corollary}
\DeclareMathOperator{\wt}{wt}
\DeclareMathOperator{\imb}{imb}
\newcommand{\ora}[1]{\overrightarrow{#1}}

\title[A union of two non-DDMO copies of $K_{3,3}$]
{A Union of Two Non-DDMO Copies of $K_{3,3}$\\ Is Difference Distance Magic Orientable}

\date{}

\subjclass[2020]{05C78, 05C20, 05C22}
\keywords{difference distance magic labeling, oriented graph, DDMOG, disconnected graph,
disjoint union, complete bipartite graph}

\begin{document}

\begin{abstract}
Aceska, Arcila-Maya, Carlson, Marr, Parnes, Ryan, Schuerger and Vasquez close their paper on
difference distance magic labelings by asking: \emph{is it possible for a union of non-DDMOGs
to be DDMO?} We answer yes. Let $\ora{D}$ be the orientation of $K_{3,3}$ with parts
$A=\{0,1,2\}$, $B=\{3,4,5\}$ and arcs
$0\!\to\!3$, $4\!\to\!0$, $5\!\to\!0$, $3\!\to\!1$, $1\!\to\!4$, $1\!\to\!5$,
$3\!\to\!2$, $2\!\to\!4$, $2\!\to\!5$.
Then $\ora{D}\sqcup\ora{D}$ carries the difference distance magic labeling
$(6,1,5,10,2,8)$, $(12,3,9,11,4,7)$ of $\{1,\dots,12\}$, while $K_{3,3}$ itself has no
orientation admitting such a labeling. So a disjoint union of two graphs, \emph{neither} of
which is difference distance magic orientable, is difference distance magic orientable.
\end{abstract}

\maketitle

\section{The question}\label{sec:q}

An \emph{oriented graph} is a digraph with no pair of opposite arcs. Following
Altman, Calzado, Freyberg, Kov\'a\v{r}, Lewis, Marr, Ross and Via~\cite{Altman}
and Aceska et al.~\cite{Source}, for a vertex $v$ of an oriented graph $\ora{G}$ on $n$
vertices and a labeling $f$ of its vertices put
\[
  \wt(v)\;=\;\sum_{u\to v} f(u)\;-\;\sum_{v\to u} f(u),
\]
the in-label sum minus the out-label sum. A \emph{difference distance magic} (DDM) labeling of
$\ora{G}$ is a bijection $f\colon V(\ora{G})\to\{1,\dots,n\}$ with $\wt(v)=0$ for
\emph{every} vertex $v$; an oriented graph possessing one is a \emph{DDMOG}, and an undirected
graph is \emph{DDMO} when at least one of its orientations is a DDMOG.

The source paper~\cite{Source} ends with a paragraph of open questions. Its third question is
unnumbered prose in the Conclusion (Section~13), in no theorem-like environment, and reads,
verbatim:

\begin{quote}
For example, is it possible for a union of non-DDMOGs to be DDMO?
\end{quote}

\noindent
Its immediate context, verbatim, is: ``\emph{Furthermore, although most of the DDMOGs in this
paper have been connected, disconnected DDMOGs are of interest as well. For example, is it
possible for a union of non-DDMOGs to be DDMO?}'' Throughout, ``union'' means disjoint union,
which is the source's own usage: its Theorem~5.1 writes a disconnected oriented graph as
$\bigcup_i\ora{G_i}$.

The question unions \emph{oriented} objects and then applies the \emph{undirected} predicate,
so it has two inequivalent readings, and we answer both. Fix $k\ge 2$.
\begin{itemize}
\item[\textbf{(A)}] (orientations fixed throughout) There are oriented graphs
$\ora{G_1},\dots,\ora{G_k}$, none a DDMOG, with $\ora{G_1}\sqcup\cdots\sqcup\ora{G_k}$ a DDMOG.
\item[\textbf{(B)}] (undirected) There are graphs $H_1,\dots,H_k$, none DDMO, with
$H_1\sqcup\cdots\sqcup H_k$ DDMO.
\end{itemize}
An instance of (B) yields one of (A), because a disjoint union has no arcs between parts, so
any DDMOG orientation of $H_1\sqcup\cdots\sqcup H_k$ restricts to orientations of the $H_i$,
each a non-DDMOG. We require $k\ge2$.

\begin{theorem}\label{thm:main}
Reading \textup{(B)}, and hence reading \textup{(A)}, holds with $k=2$: two disjoint copies of
$K_{3,3}$ form a DDMO graph on $12$ vertices, although $K_{3,3}$ is not DDMO.
\end{theorem}

No minimality of the order $12$ is claimed. The mechanism is one line: a component of a
\emph{disconnected} DDMOG on $n$ vertices receives an arbitrary $n_i$-\emph{subset} of
$\{1,\dots,n\}$, whereas being a DDMOG in its own right demands the \emph{interval}
$\{1,\dots,n_i\}$.

Since $\wt$ is defined by an equality to zero, replacing $\wt$ by $-\wt$ changes nothing
below; this matters only because~\cite{Source} reverses the usual superscript convention, so
that its $N^{+}$ is the in-neighbourhood. All statements are in the convention displayed
above.

\section{The witness on $12$ vertices}\label{sec:w12}

Let $\ora{D}$ be the orientation of $K_{3,3}$ on $A=\{0,1,2\}$, $B=\{3,4,5\}$ with the nine
arcs
\[
\begin{array}{lll}
 0\to 3, & 4\to 0, & 5\to 0,\\
 3\to 1, & 1\to 4, & 1\to 5,\\
 3\to 2, & 2\to 4, & 2\to 5.
\end{array}
\]
Each of the nine pairs in $A\times B$ occurs exactly once and no pair occurs twice, so
$\ora{D}$ is indeed an orientation of $K_{3,3}$; its in-degrees are $(2,1,1,1,2,2)$ and its
out-degrees $(1,2,2,2,1,1)$.

For this orientation the six weight equations collapse. Reading the arcs off,
\begin{align*}
 &\wt(0)=f_4+f_5-f_3, &&\wt(1)=\wt(2)=f_3-f_4-f_5,\\
 &\wt(3)=f_0-f_1-f_2, &&\wt(4)=\wt(5)=f_1+f_2-f_0,
\end{align*}
so every row of the weight matrix is $\pm(f_0-f_1-f_2)$ or $\pm(f_3-f_4-f_5)$ and
\begin{equation}\label{eq:collapse}
 \wt\equiv 0
 \iff
 f_0=f_1+f_2 \ \text{ and }\ f_3=f_4+f_5 .
\end{equation}
Equivalently: the label sets realised by $\ora{D}$ are exactly the $6$-sets of the form
$\{p,q,p+q\}\cup\{r,s,r+s\}$.

\begin{proposition}\label{prop:w12}
Let $\ora{G}_{12}=\ora{D}\sqcup\ora{D}$, and label the first copy by
$(f_0,\dots,f_5)=(6,1,5,10,2,8)$ and the second by $(12,3,9,11,4,7)$. This is a DDM labeling
of $\ora{G}_{12}$.
\end{proposition}

\begin{proof}
The two label sets $\{1,2,5,6,8,10\}$ and $\{3,4,7,9,11,12\}$ are disjoint with union
$\{1,\dots,12\}$, so the labeling is a bijection onto $\{1,\dots,12\}$. By
\eqref{eq:collapse} all twelve weights vanish as soon as
\[
 1+5=6,\qquad 2+8=10,\qquad 3+9=12,\qquad 4+7=11,
\]
which is four additions. Equivalently, the witness is the partition
$\{1,\dots,12\}=\bigl(\{1,5,6\}\cup\{2,8,10\}\bigr)\sqcup\bigl(\{3,9,12\}\cup\{4,7,11\}\bigr)$
into four triples $\{a,b,a+b\}$.
\end{proof}

\section{Neither copy is DDM-orientable}\label{sec:not}

\begin{lemma}[Altman et al.~\cite{Altman}, Theorem~5, at $m=n=3$]\label{lem:k33}
No orientation of $K_{3,3}$ has a DDM labeling. Hence $K_{3,3}$ is not DDMO and every
orientation of $K_{3,3}$ is a non-DDMOG.
\end{lemma}

This is not new: \cite[Theorem~5]{Altman} states that $K_{m,n}$ with $m,n\ge3$ has an
orientation permitting a DDM labeling if and only if $(m+n)(m+n+1)=4s$ for some
$s\in\mathbb{N}$, and $(3+3)(3+3+1)=42$ is not a multiple of $4$. That paper is the first item
in the bibliography of~\cite{Source}. We reproduce the argument, which is the necessity half of
theirs, only so that the witness is self-contained.

\begin{proof}
Let $\ora{K}$ be any orientation of $K_{3,3}$ and let $f$ be any labeling with all weights
zero. Summing the weights over $A$: every arc of $\ora{K}$ joins $A$ to $B$, so no $f(a)$ with
$a\in A$ occurs in $\sum_{a\in A}\wt(a)$, while $f(b)$ occurs with coefficient
$\operatorname{outdeg}(b)-\operatorname{indeg}(b)=2\operatorname{outdeg}(b)-3$, which is odd
because $\deg(b)=3$. Hence $0=\sum_{a\in A}\wt(a)\equiv\sum_{b\in B}f(b)\pmod 2$, so
$S_B:=\sum_{b\in B}f(b)$ is even, and symmetrically $S_A$ is even. But a DDM labeling has
$S_A+S_B=1+2+\cdots+6=21$, which is odd; a contradiction.
\end{proof}

Lemma~\ref{lem:k33} with Proposition~\ref{prop:w12} proves Theorem~\ref{thm:main} and answers
the question in reading~(B). The
union defeats the parity obstruction because each copy receives a $6$-\emph{subset} of
$\{1,\dots,12\}$ of even sum ($32$ and $46$), whereas the \emph{interval} $\{1,\dots,6\}$ has
odd sum $21$.

\section{Relation to Theorem~5.1 of the source}\label{sec:src}

The witness lies outside the hypothesis of Theorem~5.1 of~\cite{Source}, which asserts that
$\ora{G}=\bigcup_{i=1}^{k}\ora{G_i}$ is a DDMOG if $\ora{G_1}$ is a DDMOG and
$\ora{G_2},\dots,\ora{G_k}$ are DDMOGs with $\imb(\ora{G_i})=0$: here no component is a DDMOG
at all, by Lemma~\ref{lem:k33}. That theorem is stated as a sufficient condition only, and the
source itself declines the case the witness inhabits --- the prose introducing it reads
``\emph{However, if a union of DDMOGs happens to contain more than one component that is not
balanced, it is not immediately clear if a DDM labeling will exist for this union}''.

\section{Verification, and what is not settled}\label{sec:ver}

\noindent\textbf{By hand.} Everything above is hand-sized and needs no machine: the
$12$-vertex witness is four additions, and Lemma~\ref{lem:k33} is a parity count.

\medskip\noindent\textbf{By machine.} The accompanying program \texttt{verify.py}
(Python~3.9+, standard library only), in exact integer and rational arithmetic, reads the
objects printed above and re-checks them, and confirms Lemma~\ref{lem:k33} by two independent
routes over all $512$ orientations of $K_{3,3}$, which agree. Its recorded transcript also
carries checks on further material --- minimum orders for witnesses, order-$5$ enumerations, and
a parity condition for $k$ copies of $\ora{D}$ --- none of which is claimed in this note and
none of which anything above depends on; where that transcript names a lemma or a remark of the
paper for a step it did not re-run, no such numbered statement appears above.

\medskip\noindent\textbf{What is not settled.}
\begin{itemize}
\item No theorem describing which orders admit witnesses is proved here, and no witnesses are
counted: Theorem~\ref{thm:main} exhibits one witness, on $12$ vertices, and nothing here rules
out a witness of smaller order.
\item The published census of $6$-vertex DDMOGs in~\cite{Altman} --- ``\emph{there are 22
unique 6-vertex oriented graphs on two underlying non-isomorphic graphs with DDM labelings}''
--- is \emph{not} reproduced here and no agreement with it is claimed. The only order-$6$
object this note uses is $K_{3,3}$, whose $512$ orientations are swept exhaustively. Likewise
the $7$- and $8$-vertex counts of~\cite{Altman} play no role.
\item The ingredient the argument leans on is published, not new: Lemma~\ref{lem:k33} is
\cite[Theorem~5]{Altman} at $m=n=3$, and the proof reproduced here is the necessity half of
theirs. What is claimed here is the answer to the question and the $12$-vertex witness.
\item Besides \cite{Source} and~\cite{Altman}, the only further paper we located on difference
distance magic labelings is~\cite{Ornaments}, the one paper we found citing the source; its
constructions all adjoin a neutralising edge set and so produce connected DDMOGs, and none of
its own open questions restates or answers the one answered here. The words ``disjoint'',
``union'', ``component'' and ``disconnect'' do not occur in~\cite{Altman}.
\item The other two open questions of the same paragraph of~\cite{Source} --- which graphs are
DDMO, and what properties distinguish DDMOGs from non-DDMOGs --- remain open.
\item Splitting an interval of labels into blocks of prescribed sum is the standard device of
the undirected and group distance-magic literature (see e.g.~\cite{CSP}); the mechanism of the
$12$-vertex witness is recognisably that device, transplanted to the oriented setting.
\end{itemize}

\begin{thebibliography}{9}

\bibitem{Source}
R.~Aceska, C.~Arcila-Maya, R.~Carlson, A.~Marr, P.~Parnes, J.~Ryan, H.~Schuerger, and
A.~Vasquez,
\emph{New results on difference distance magic labelings},
arXiv:2601.20492v1, 2026.
\href{https://arxiv.org/abs/2601.20492}{arXiv:2601.20492};
\href{https://doi.org/10.48550/arXiv.2601.20492}{doi:10.48550/arXiv.2601.20492}.
The question answered here is unnumbered prose in the Conclusion (Section~13) of that version,
and is quoted in full in Section~\ref{sec:q}.

\bibitem{Altman}
K.~Altman, L.~Calzado, B.~Freyberg, P.~Kov\'a\v{r}, E.~Lewis, A.~Marr, L.~Ross, and R.~Via,
\emph{Difference distance magic oriented graphs},
Res. Math. Sci. \textbf{11} (2024), no.~4, art.~67.
\href{https://doi.org/10.1007/s40687-024-00482-7}{doi:10.1007/s40687-024-00482-7}.

\bibitem{Ornaments}
R.~Aceska, T.~Lort, and A.~Ripperger,
\emph{Ornaments and difference distance magic oriented graphs},
arXiv:2606.16001v3, 2026.
\href{https://arxiv.org/abs/2606.16001}{arXiv:2606.16001}.
Every DDMOG constructed there is connected, so none of its results bears on the question
answered here.

\bibitem{CSP}
B.~Freyberg,
\emph{On constant sum partitions and applications to distance magic-type graphs},
AKCE Int. J. Graphs Comb. \textbf{17} (2020), no.~1, 373--379.
\href{https://doi.org/10.1016/j.akcej.2019.02.003}{doi:10.1016/j.akcej.2019.02.003}.

\end{thebibliography}

\end{document}
